%%%----------------------------------------------------------------------------------------
%	SECTION TITLE
%----------------------------------------------------------------------------------------
\vspace{-3mm}
\cvsection{Experience}

%----------------------------------------------------------------------------------------
%	SECTION CONTENT
%----------------------------------------------------------------------------------------

\begin{cventries}

%------------------------------------------------

\cventry
{Staff Engineer II (Engienering Physicist II)} % Job title
{SLAC National Laboratory} % Organization
{Palo Alto, CA} % Location
{April. 2022 - PRESENT} % Date(s)
{ % Description(s) of tasks/responsibilities
\begin{cvitems}
%\item {Developed analysis workflows to correlate bunch length measurements (BLEN) with L2 phase and XTCAV diagnostics}
\item {Led the recommissioning of the bunch length diagnostic system after six years of dormancy, restoring a critical beam diagnostic and enabling reliable, high-resolution measurements for accelerator operations and research.}
\item {Built MATLAB-based tools to process and visualize beam diagnostics data across accelerator sectors}
\item {Contributed to accelerator operations and troubleshooting, supporting beam delivery and system optimization}
\item {Led the restoration and recommissioning of the FACET positron beamline, a dormant accelerator section inactive since 2016, returning the system to full operational status.}
\item {Directed cross-disciplinary engineering efforts spanning vacuum systems, electrical infrastructure, and water cooling systems to re-establish beamline functionality and ensure safe, reliable operation.}
\item {Redesigned and implemented simplified electrical wiring architectures for legacy FACET magnet systems, improving system clarity, maintainability, and operational efficiency.}
\item {Engineered electrical configurations to enable clear and enforceable Lockout-Tagout (LOTO) conditions, significantly improving safety and accessibility for maintenance and operations personnel.}
\end{cvitems}
}
\vspace{3mm}
%------------------------------------------------    

\cventry
{Engineering Physiscist I, US-HL-AUP Superconducting R\&D Engineer \& Electrical Quality Assurance } % Job title
{Fermilab} % Organization
{Batavia, IL} % Location
{April 2015 - March 2022} % Date(s)
{% Description(s) of tasks/responsibilities
\begin{cvitems}
% tryint o say that I led the daily active floor fabrication efforts for super cold magnet coils. I was in charge of daily priorities baseon on expected deadlines for magnet shipping. I kept track of the rates at which we were succesfully outputting coils, and collaborated with technicians regarding 
\item {Managed daily fabrication workflow for superconducting magnet coils, optimizing production rates and addressing process inefficiencies in collaboration with technicians to meet delivery targets.}
\item {Led electrical quality control for Nb3Sn superconducting magnet systems in the US HiLumi-LHC AUP, improving fabrication reliability and reducing defect rates during early-stage production.}
\item {Lead successful design and fabrication of R\&D test Coils used used to investigate failure modes of in epoxy injected superconducting cables.}
\item {Engaged in electromechanical failure analyses of prototype and production model Nb3Sn magnet coils for the AUP Project}
\item {Initiated and lead a review devoted to enhancing and updating electrical design parameters applied to AUP magnet coils during early fabrication stages}
\item {Designed and implemented a robust nitrogen (N$_2$) distribution system for furnace operations to displace oxygen during coil processing, improving reliability and reducing technician intervention.}
\item {Supervised electro-mechanical technicians through special R\&D project for cold-magnet coil fabrication efforts}
\item {Worked closely with physicists to optimise and improve the heat treatment process for the AUP Coils to achieve the desired superconducting performace goals based on laboratory testing and analysis} 
\item {Executed final electrical qualification of AUP superconducting coils prior to shipment to CERN, verifying compliance with specifications and ensuring readiness for integration.}
\end{cvitems}
}
\vspace{3mm}
%------------------------------------------------

\cventry
{Accelerator Systems \& Operations} % Job title
{Fermilab} % Organization
{Batavia, IL} % Location
{April 2014 - Feb. 2013} % Date(s)
{ % Description(s) of tasks/responsibilities
\begin{cvitems}
\item {Coordinated accelerator operations to initiate, optimize, and maintain beamline performance in support of laboratory and university-based research programs.}
\item {Operated, tuned, and maintained the H$^-$ plasma source to ensure stable beam delivery across the accelerator complex.}
\item {Contributed to the design and implementation of an air diversion system for the NuMI horn power stripline, improving thermal management and operational reliability for high-power beam delivery.}
\end{cvitems} 
}
\vspace{6mm}
%------------------------------------------------
\cventry
{Graduate Research Assistant} % Job title
{Georgia Institute of Technology} % Organization
{Atlanta, GA} % Location
{Dec. 2013 - April. 2014} % Date(s)
{ % Description(s) of tasks/responsibilities
\begin{cvitems}
\item {Participated in an in-depth study to understand the interaction of a Large plasma boundaries and their associated wall interactions}
\item {Studied plume divergence behavior of the P5 - 5 Kw hall thruster}
\item {Assisted in the design of the glass injections tube for a helicon ion thruster}
\item {Aided in the fabrication of several hallow cathodes and plasma probes for testing}
\end{cvitems}
}
\vspace{6mm}
%------------------------------------------------
\cventry
{Active Aerolasticity and Structures Lab - Undergraduate Research Assistant } % Job title
{University Of Michigan} % Organization
{Ann Arbor, MI} % Location
{Dec. 2012 - April. 2013} % Date(s)
{ % Description(s) of tasks/responsibilities
\begin{cvitems}
\item Conducted aeroelasticity research on highly flexible aircraft, supporting design and testing of UAV systems.
\item Designed and validated a ventral fin using CFD and wind tunnel testing (2'×2' and 5'×7') to improve yaw stability and control authority.
\item Developed a sensor-based system using gyroscopes to model and measure aircraft structural deformation in flight.
\item Fabricated and maintained structural components for the X-HALE UAV, supporting iterative design improvements.
\item Contributed to full-scale assembly of an Aerial Test Vehicle (ATV), including composite wing layup and manufacturing.
\end{cvitems}
}
%------------------------------------------------
%------------------------------------------------
\end{cventries}